% Capítulo extra, ainda não escrito no percurso principal.
% Este ficheiro é deliberadamente independente e não é incluído por main.tex.
% Descreve uma proposta de protocolo; não descreve o consenso actual de Monero.

\chapter{O nó é a linhagem}
\label{chapter:unwritten-node-identity}

\begin{quote}
\lq\lq Nunca digo a minha chave. De tempos a tempos, demonstro apenas que ainda a
conheço.\rq\rq
\end{quote}

Uma máquina não atravessa a rede. A rede vê mensagens, tempos, blocos e provas.
Se quisermos ligar essas observações a uma reputação, não precisamos de afirmar
que vieram sempre da mesma caixa de plástico, da mesma placa ou sequer do mesmo
continente. Podemos dar ao nó uma definição mais estreita e, por isso mesmo,
mais útil: o nó é uma linhagem de chaves que o protocolo reconhece como uma só
identidade.

Esta distinção muda a pergunta. Já não tentamos provar que um sistema operativo
particular está agora a correr numa máquina física particular. Perguntamos se o
interlocutor conhece a chave corrente da linhagem e se cada mudança de chave foi
autorizada pela anterior. Uma prova de Schnorr responde à primeira pergunta; uma
cadeia de rotações assinadas responde à segunda.

O resultado pode prender histórico de mineração, disponibilidade, recompensas e
penalizações a um nó sem revelar a sua chave privada. Não é uma atestação remota
do estado de execução. É uma identidade criptográfica com memória.

\section{De uma imagem de disco a uma chave escondida}

Seja \(B\) uma sequência canónica de bytes: por exemplo, todos os sectores de
uma partição imutável, por uma ordem e com um comprimento precisamente
definidos. Calculamos primeiro
\[
    d=\mathcal H_{\rm disk}(\mathsf{partition\mbox{-}v1},B).
\]
O rótulo de domínio impede que o mesmo resumo seja interpretado por acidente
como um valor criado para outro protocolo. Para obter um escalar do grupo,
podemos calcular
\[
    x=\mathcal H_n(\mathsf{node\mbox{-}key\mbox{-}v1},d,\rho),
    \qquad X=xG,
\]
onde \(\rho\) é entropia privada e imprevisível própria daquele nó.

O ponto \(X\) pode ser publicado; \(x\) e \(d\) não precisam de sair do nó. Sob
a hipótese do logaritmo discreto, observar \(X=xG\) não permite recuperar
\(x\). As provas de Schnorr posteriores também não o revelam.

A entropia \(\rho\) não é um enfeite. Se duas máquinas usarem exactamente a
mesma imagem pública e definirem simplesmente \(x=\mathcal H_n(d)\), ambas obtêm
a mesma chave. Além disso, qualquer pessoa que possua uma imagem candidata pode
calcular o seu \(d\), calcular \(\mathcal H_n(d)G\) e compará-lo com \(X\). O
resumo não se torna secreto só porque nunca é enviado. A segurança da chave
exige material privado com entropia suficiente; nesse sentido, \(\rho\) é já uma
chave, ainda que a imagem participe na sua derivação.

Há outra fronteira que convém traçar cedo. Um provador sem ambiente de execução
fiável pode guardar \(x\), guardar uma cópia limpa de \(B\), ou calcular a prova
enquanto executa outro programa. A construção demonstra conhecimento de um
escalar derivado dos bytes; não demonstra que esses bytes são precisamente o
código que o processador executa naquele instante.

\section{Bater à porta com Schnorr}

O nó regista \(X=xG\). Ao abrir uma ligação, o outro participante envia um
\emph{nonce} fresco \(q\). O contexto autenticado deve conter mais do que esse
nonce. Um exemplo é
\[
 m=(\mathsf{node\mbox{-}link\mbox{-}v1},\mathsf{chainID},q,
       \mathsf{keys}_{\rm session},\mathsf{tip}),
\]
onde \(\mathsf{keys}_{\rm session}\) liga a prova à negociação criptográfica da
ligação e \(\mathsf{tip}\) identifica o estado da cadeia que os participantes
pensam estar a discutir.

O nó escolhe \(\alpha\in_R\mathbb Z_l\), calcula
\begin{align*}
    A&=\alpha G,\\
    c&=\mathcal H_n(\mathsf{node\mbox{-}schnorr\mbox{-}v1},m,X,A),\\
    r&=(\alpha-cx)\bmod l
\end{align*}
e envia \((A,r)\). O interlocutor recalcula \(c\) e aceita quando
\[
    A=rG+cX.
\]
Como sempre, reutilizar \(\alpha\) com desafios diferentes revela \(x\). Uma
implementação deve ainda validar as codificações e os pontos, separar domínios
e prender ao desafio todos os campos cujo significado se quer autenticar.

Duas sessões válidas em momentos diferentes, verificadas contra o mesmo \(X\),
demonstram conhecimento do mesmo \(x\) em cada um desses momentos. Não
demonstram que \(x\) permaneceu continuamente em memória entre as sessões, que
nunca foi copiado ou que só uma máquina o conhece. Para a identidade do nó, a
declaração necessária pode ser precisamente a mais pequena: quem controla a
chave corrente fala em nome do nó.

\section{A identidade sobrevive à chave}

Uma chave que nunca muda acumula risco juntamente com reputação. Para permitir
rotação, começamos por uma chave \(X_0=x_0G\) e damos à linhagem um nome estável
\[
    N=\mathcal H_{\rm node}(\mathsf{nodeID\mbox{-}v1},X_0).
\]
Na geração \(i\), o nó escolhe independentemente
\(x_{i+1}\in_R\mathbb Z_l\), calcula \(X_{i+1}=x_{i+1}G\) e constrói a
declaração
\[
 T_i=(\mathsf{rotate\mbox{-}v1},\mathsf{chainID},N,i,
          X_i,X_{i+1},h_{\rm activation}).
\]
A chave corrente autoriza a sucessora:
\[
    \sigma_i=\operatorname{Sign}_{x_i}(T_i).
\]
A nova chave deve também assinar \(T_i\), ou um resumo canónico dele. Esta
contra-assinatura demonstra que o autor da rotação não apontou a linhagem para
uma chave inutilizável ou desconhecida.

O estado reconhecido pelo protocolo toma então a forma
\[
 X_0\xrightarrow{\,T_0\,}X_1
    \xrightarrow{\,T_1\,}X_2
    \xrightarrow{\,T_2\,}\cdots .
\]
Uma regra completa de rotação precisa de decidir, pelo menos, o seguinte:
\begin{itemize}
    \item a geração seguinte é exactamente \(i+1\);
    \item só a chave corrente pode autorizar a sucessora;
    \item a chave nova prova que existe antes de a rotação ser aceite;
    \item depois de \(h_{\rm activation}\), mensagens da chave antiga são
          rejeitadas;
    \item duas sucessoras assinadas para a mesma geração constituem um conflito
          publicamente demonstrável;
    \item a regra de consenso determina qual rotação se torna canónica e em
          que momento fica final.
\end{itemize}

Uma rotação não deve ser apenas \(x_{i+1}=\mathcal H_n(x_i)\). Quem roubar
\(x_i\) conseguiria calcular todas as chaves futuras dessa cadeia. Gerar a
chave seguinte independentemente e entregar apenas a sua chave pública à
anterior evita essa consequência simples, embora o roubo da chave corrente
continue grave.

\section{Reputação: a memória da rede}

Guardemos a reputação sob \(N\), não sob cada \(X_i\):
\[
    \mathsf{Rep}[N]
       =F(\mathsf{trabalho},\mathsf{disponibilidade},
           \mathsf{antiguidade},\mathsf{faltas}).
\]
Uma chave nova e sem certificado começa uma linhagem nova, portanto começa sem
o histórico de \(N\). Uma chave nova acompanhada pela rotação válida herda a
mesma identidade e o protocolo continua a actualizar a mesma entrada de
reputação.

É aqui que a construção ganha força económica. Um actor pode criar milhares
de pares de chaves, mas nenhum deles herda gratuitamente o trabalho ou a boa
conduta atribuídos à linhagem antiga. Para falar com a reputação de \(N\), tem
de conhecer a chave corrente de \(N\) ou apresentar uma transição que ela
autorizou.

Isto não obriga um atacante a executar o programa honesto. Obriga-o a obedecer
à máquina de estados da identidade se quiser conservar aquela reputação. Pode
abandoná-la e começar de novo; não pode transportar a reputação para uma chave
arbitrária sem a autorização exigida.

A identidade, por si só, também não resolve o problema de Sybil. Apenas dá à
rede gavetas persistentes onde guardar memória. O custo de encher uma dessas
gavetas com boa reputação tem de vir do mecanismo escolhido: trabalho válido,
tempo, recursos escassos, garantias sujeitas a perda ou outro sinal que não
possa ser fabricado em massa sem custo.

\section{Prender o trabalho de mineração ao nó}

Uma acção de mineração atribuída à linhagem deve incluir \(N\), a geração
\(i\) e \(X_i\) no contexto autenticado. Esquematicamente,
\[
 M=(\mathsf{mining\mbox{-}v1},\mathsf{chainID},N,i,X_i,h,
       \mathsf{previousBlock},\mathsf{workObject}).
\]
O nó apresenta o trabalho exigido pelo protocolo e uma assinatura válida de
\(M\) sob \(X_i\). Incluir a identidade no próprio objecto sobre o qual o
trabalho é calculado pode impedir que uma prova observada seja simplesmente
reatribuída a outro \(N\). A assinatura Schnorr não cria trabalho e não
substitui a regra de mineração; autentica apenas a atribuição daquele trabalho
à linhagem.

Nonces de ligação, alturas, identificadores da cadeia e o bloco anterior têm
papéis diferentes. O nonce impede repetir uma autenticação de sessão; a altura
e a geração escolhem a chave válida; \(\mathsf{chainID}\) impede levar uma
assinatura para outra rede; o bloco anterior prende o trabalho à história em
que foi produzido. O protocolo deve codificar todos eles sem ambiguidades.

\section{O roubo da chave corrente}

Se um atacante obtiver \(x_i\), passa a poder falar como \(N\), assinar trabalho
e autorizar \(X_{i+1}\). A rotação limita a exposição de chaves antigas, mas
não distingue duas pessoas que conhecem a mesma chave corrente. Sem hardware
fiável, a rede não consegue provar que só existe uma cópia nem que \(x_i\) foi
apagado depois da rotação.

Uma defesa possível separa a identidade de reputação da chave operacional.
Seja
\[
    Y=yG
\]
uma chave-raiz guardada fora do mineiro em linha. A reputação fica ancorada em
\(Y\), enquanto \(Y\) certifica chaves \(X_i\) de curta duração. A rotação
normal pode exigir a chave operacional antiga e a nova; uma recuperação ou uma
mudança excepcional pode exigir \(y\), talvez com atraso público. Assim, roubar
uma chave de mineração não entrega necessariamente toda a linhagem para sempre.

Outras políticas podem usar múltiplas chaves de recuperação, atrasos de
activação, períodos de contestação ou penalizações por rotações incompatíveis.
São escolhas de disponibilidade e governação, não consequências automáticas de
Schnorr. Uma raiz fora de linha também troca a elegante afirmação
\lq\lq a imagem é a chave\rq\rq por uma arquitectura mais honesta: a imagem pode contribuir para a
chave operacional, mas a continuidade de reputação repousa num segredo próprio.

\section{O que ficou demonstrado}

Podemos agora separar quatro afirmações que à primeira vista se confundem:
\begin{enumerate}
    \item Uma prova válida sob \(X_i\) demonstra conhecimento de \(x_i\) no
          momento da prova.
    \item Uma rotação válida demonstra que \(x_i\) autorizou \(X_{i+1}\) e
          que alguém demonstrou possuir \(x_{i+1}\).
    \item A cadeia de rotações permite ao protocolo tratar todas essas chaves
          como uma identidade \(N\) e conservar a sua reputação.
    \item Nada disto, sem uma raiz externa de confiança, demonstra qual código
          está a correr, qual máquina o executa, quantas cópias da chave existem
          ou se uma chave antiga foi realmente apagada.
\end{enumerate}

A quarta frase não desfaz as três primeiras. Define-lhes a fronteira. Para uma
rede cuja unidade social e económica é o nó criptográfico, e não uma caixa
física, essa fronteira pode ser exactamente a desejada. O nó não é a máquina.
O nó é a história verificável de quem teve autoridade para falar a seguir.
